% =====================================================================
%  THE LEDGER FRAMEWORK — Market Data Manifesto, Version 1.1
%  Governing principles for market data. Subordinate to the
%  Architectural Constitution (ledger_manifesto_v1_4.tex).
%  1.1 amends the certified 1.0 (fixed) with three embedded mandates:
%  (1) corporate actions as first-class events — frames and the market
%  data operator algebra (MD-13); (2) a new governing principle,
%  dispute-readiness (MD-14); (3) model association and price-space
%  validation (MD-15), with MD-1/MD-6/MD-9 extended. MD-1..MD-12 keep
%  their numbers.
%  Drafting lead: GATHERAL. Rigor-check: FORMALIS. Review: THORP,
%  KLEPPMANN, TALEB. Certification: CONCORDIA.
%  Prior art (source material, not a constraint): the Market Data
%  Calibration Manifesto (calibration_manifesto.tex).
%  Build: pdflatex (twice). Standard TeX Live; lmodern and microtype
%  used when available and skipped otherwise.
% =====================================================================
\documentclass[11pt,a4paper]{article}

\IfFileExists{lmodern.sty}{%
  \usepackage[T1]{fontenc}\usepackage{lmodern}%
  \IfFileExists{microtype.sty}{\usepackage{microtype}}{}%
}{}
\usepackage[margin=2.5cm,marginparwidth=2.0cm,marginparsep=2mm]{geometry}
\usepackage{parskip}
\usepackage[hidelinks]{hyperref}

% --- article addressing: margin identifiers, mirroring the Constitution ---
\IfFileExists{xcolor.sty}{\usepackage{xcolor}}{}
\ifdefined\textcolor\else\def\textcolor#1#2{#2}\fi
\newcommand{\mlabel}[1]{\leavevmode\hypertarget{md-#1}{}%
  \marginpar{\raggedright\scriptsize\ttfamily\textcolor{gray}{#1}}}

\setcounter{secnumdepth}{0}
\providecommand{\tightlist}{\setlength{\itemsep}{2pt}\setlength{\parskip}{0pt}}

% traceability tag at the close of each article
\newcommand{\trace}[1]{\par\smallskip{\small\itshape (#1)}}

\begin{document}

\begin{center}
  {\LARGE\bfseries THE LEDGER FRAMEWORK}\\[0.9em]
  {\large\bfseries Market Data Manifesto}\\[0.6em]
  {\normalsize Governing Principles for Market Data ---
   Subordinate to the Architectural Constitution}\\[1.4em]
  Version 1.1 \quad---\quad amends the certified 1.0 \quad---\quad July 2026
\end{center}

\vspace{0.6em}
\begin{center}\rule{0.35\linewidth}{0.4pt}\end{center}
\vspace{0.6em}

\begin{center}\itshape
The record holds observations; every price beyond them is derived from them.
\end{center}

\vspace{0.4em}

\begin{center}\textbf{Abstract}\end{center}
\begin{quote}
The Ledger records observations, not prices. A market data fact enters
the record exactly once, as a moveless observation admitted through the
one door; every curve, surface, correlation, and calibration is a
deterministic function of recorded inputs, of two kinds the document
keeps apart. A \emph{projection} computes over the record and rebuilds
from it alone; a model's output \emph{re-enters as a recorded
observation} (C-14.9, C-14.15). Reading back any recorded value is
unconditional; re-deriving a model's output is what needs the model. A
value is meaningful only in a \emph{frame} and at a time coordinate ---
corporate actions are first-class events that change frames; every datum
is bound to a model by recorded lineage and validated in price space;
and every datum is dispute-ready, settled by replay. This manifesto
states the principles that make append-only, event-driven, projected
market data the only sensible way to hold it. Each principle is an
\emph{article} derived from six of the Constitution's eight commitments
--- auditability, reproducibility, time travel, correctness, order,
simulability --- and a seventh it \emph{derives} from the first two,
dispute-readiness; each rests on the Constitution clause that carries it.
The document is rated against the Constitution, never the other way round;
where the two conflict, the Constitution prevails.
\end{quote}

\vspace{0.4em}
\begin{center}\rule{0.35\linewidth}{0.4pt}\end{center}
\vspace{0.6em}

\section{1. Market Data in the Ledger's Terms}

A market data fact is an \textbf{observation}: a recorded statement,
attributed to a source, that a named observable --- a price, a rate, a
fixing, an index level --- took a value with respect to a world-moment.
The observation is a fact about what a source said, not the truth about
what a thing is worth; any ``true'' price behind it is a modelling posit
that lives in a model, never a claim the record makes.

The source of an observation sits outside the trust boundary; the Events
Executor captures and records the arrival (C-5.4). It enters the immutable
log only as a moveless \emph{observation-recording transaction} ---
moveless because it records data and moves no balances --- admitted
through the single Transaction Executor and folded into a \emph{home}, its
resting place in the record from which contracts read (C-4.8). (To
\emph{fold} the log is to replay it in order, accumulating state; the
current state is the initial state plus the fold of all observations.)
There is no second door for market data: the one record and the single
writer extend to every observation. A kind of
market data is declared before any observation of it crosses --- its
fields, the frame it is delivered in, the market data operators that
adjust it across corporate actions (C-9.2), and the grain of the
cause-derived identifier --- an identifier computed
from what caused an arrival, so a redelivery of the same cause collides
with it --- that decides when two arrivals are the same fact, all
registered on the record first, so an arrival of an undeclared kind is
quarantined, never guessed (C-4.12).

\textbf{Vocabulary, mapped once.} This manifesto uses the Constitution's
fixed names and coins no synonym. Its prior art, the Market Data
Calibration Manifesto, speaks a different dialect, mapped here once and
not carried further: an \emph{Oracle} is an external source (no
``oracle'' primitive, MD-3); an \emph{attestation} is an observation
(C-4.8); \emph{event time} is execution time, \emph{knowledge time}
splits into monitor and door time (C-2.7); a \emph{posterior} is a derived
object (MD-6); a \emph{model configuration} is the recorded, versioned
terms of a derived object (C-4.4, C-7.2). The Constitution's own
term \emph{the market data operator} names one thing only --- the
corporate-action transform of C-9.2 --- and is never used for a source;
version 1.1 uses that operator as the map between \emph{frames} (MD-13).
\emph{Frame} builds on the Constitution's own ``post-event frame''
(C-9.2), which this manifesto promotes to a defined coordinate system ---
the one term it fixes that the Constitution names only informally. No
synonym for the
operator is introduced.

\section{2. The Articles}

Each article derives from one or more of the six commitments and rests
on a clause of the Constitution. The identifiers MD-\emph{n} are
append-only: never reused, never renumbered, so later documents may cite
them.

\mlabel{MD-1}\textbf{MD-1. The observation is the atom of market data.}
A market data fact is an observation, and an observation is recorded
exactly once. It enters the immutable log as a moveless
observation-recording transaction, admitted through the single writer,
and folded into a home. ``Exactly once'' is a mechanism, not a wish: each
arrival carries its cause-derived identifier, and an arrival identical
under it is \emph{absorbed} as a duplicate, not recorded again (C-2.7,
C-12.3, C-14.3). Absorption is not a silent drop: no new value is
recorded, but the redelivery's arrival provenance is retained, so a
healthy absorb is distinguishable from a feed gone quiet. Three arrivals that look alike must be told
apart, and the identifier's grain --- a registered property of the data
kind (\S1) --- is what tells them: the same print \emph{redelivered} is
absorbed; a vendor \emph{correction} is a distinct, later observation
that names the value it corrects (MD-10); a genuinely \emph{new} print at
the same instant is a second observation, recorded in its own right. A
grain too \emph{fine} double-counts a redelivery; a grain too \emph{coarse}
absorbs a genuine second print and loses it --- and that the registered
grain is correct is a trust assumption, named beside TA-ARRIVAL. The
double-count is loud, the extra observation on the record to be found;
the over-coarse loss is \emph{silent}, and where two distinct prints are
identical under every field a feed sends, no grain can separate them.
This one residual silent loss the manifesto does not pretend to prevent;
it is caught where TA-ARRIVAL's gaps are --- at the perimeter, reconciling
arrival counts against recorded observations. Subject to that, nothing
about market data is written twice: the value a source stated is its sole
canonical copy, computed from wherever it is mentioned, never stored
beside it.
\trace{Auditability, Correctness; C-1.4, C-4.8, C-2.7, C-12.3, C-14.3.}

\mlabel{MD-2}\textbf{MD-2. Capture precedes judgement.} The record
admits an observation before anything decides whether to believe it. A
stale, fat-fingered, or off-consensus quote is well-formed: it enters
through the normal door as an observation-recording transaction (C-4.8),
and judgement is deferred because economic correctness is never gated at
admission --- it is detected after the fact by recomputation (C-13.2,
C-13.3). A quote that is not even processable --- a corrupted payload,
an unresolvable reference --- is still not lost: it is captured with the
provenance of its arrival and stands as a visible open item (C-4.12).
Either way the Ledger captures, then classifies; it never refuses to
record in order to reject. Whether an observation is used, down-weighted,
or ignored is a separate and later question, answered by a derived
object over the record, never by the door. A ``rejected'' observation is
one a calibration assigned no weight --- present, inspectable,
reclassifiable later --- not one the record never held.
Capture-then-classify is a \emph{principle} whose guarantee runs
\emph{from arrival}: whether a pipe can hold the open's firehose is a
capacity question outside this document (TA-ARRIVAL), but the principle
forbids using capacity as a licence to drop silently --- what the boundary
cannot admit is a named, visible residual; what never reaches it is the
perimeter's concern (C-4.12).
\trace{Auditability, Order, Correctness; C-4.8, C-13.2, C-13.3, C-4.12.}

\mlabel{MD-3}\textbf{MD-3. The source is untrusted; its delivery is a
trust assumption.} The origin of an observation --- an exchange, a
broker, a fixing provider, an on-chain feed --- is machinery beyond the
system's trust boundary, an external source (C-5.1, C-5.8), and the
framework grants no source the standing of truth. That an observation
reaches the boundary at all is a trust assumption, reconciled at the
perimeter, never a fact the fold may assume (C-4.12). The framework has no
privileged ``oracle'' component: a source is where an observation came
from, recorded as provenance, and nothing more.
\trace{Auditability, Correctness; C-5.1, C-5.4, C-5.8, C-4.12.}

\mlabel{MD-4}\textbf{MD-4. Every observation bears three times, and as-of
and as-at are two of them.} By the Constitution an event carries an
execution time --- the world-moment its source asserts it happened, the
time a court would enforce, contestable only in the world and corrected
only by a later event; a monitor time --- when the boundary observed it;
and a door time --- when the single writer admitted it (C-2.7). The
\emph{as-of} coordinate --- the world-moment a value concerns --- is its
execution time; the \emph{as-at} coordinate --- when the record came to
hold it --- is its door time; monitor time is provenance, splitting
lateness into the world's delay and ours. Because both are recorded, the
system computes both honest answers at any moment: what was known as of a
past date, and what is known now about it (C-12.1). A value carries one
coordinate more: the corporate-action frame it is expressed in (MD-13).
\trace{Time travel, Order; C-2.7, C-12.1.}

\mlabel{MD-5}\textbf{MD-5. Order is meaning, and the authoritative fold
is execution-ordered.} Observations are events, and their sequence is
load-bearing: a derived object depends on which observations preceded
it --- a surface fitted before a shock is not the surface fitted after
it. Meaning lives in execution order. An observation that arrives late,
whose execution time precedes observations already folded, takes its
place among them; every \textbf{projection} downstream refolds to include
it (C-2.7), while a \textbf{re-entered observation} whose input window it
lands in cannot re-fold itself --- the ledger runs no model (C-14.9) ---
and is instead flagged stale for re-derivation (MD-8). Neither is silent:
the insertion and every changed or staled value are flagged, and the
difference is published as a named explain item --- a line in the
profit-and-loss explain (C-12.6). The as-known view --- what the book
showed with only the observations it then held --- survives unrewritten
beside the corrected one.
\trace{Order, Time travel; C-2.7, C-12.6, C-14.9.}

\mlabel{MD-6}\textbf{MD-6. Every derived object is a deterministic
function of recorded inputs.} A curve, a surface, a correlation, a
calibration is a function of what the record holds --- the observations
and any derived objects it consumes, each pinned to its version and
\emph{cut} (the as-at boundary fixing which observations were in force;
MD-12), the recipe that defines it, and, where the recipe draws on
randomness, a recorded seed. Two kinds of
derived object must be told apart, and the Constitution names both.

Where the recipe computes \emph{over the record} --- a discount factor
read off recorded quotes, an adjusted value from the market data operator
--- the object is a \textbf{projection}: it stores nothing and rebuilds
from the record alone, by any party, today or at a past date (C-2.2,
C-3.2, C-14.11). Where the recipe \emph{runs a model} --- a filter, an
optimiser, a fitted surface --- the ledger does not run it (C-14.9); the
model is evaluated outside the fold and its output re-enters through the
one door as a new \textbf{observation} (C-14.15). That number is stored
like any observation, and two senses of recovering it must not be
confused: \emph{reading it back} is always possible from the record
alone; \emph{re-deriving} it --- showing it follows from its inputs ---
needs the retained model and the numerical environment it ran in, a
governance matter outside this scope. A re-entered observation is read
back unconditionally, re-derived only with the model.

Either kind carries, as recorded provenance, a \textbf{complete lineage}:
the resolved input cut --- every observation and derived object it
consumed, each pinned to version and cut; the corporate-action events
whose frame the operator applied (MD-13); for a model output the model
version; and for any datum used in valuation the model it is bound to
(MD-15). Lineage reaches through the whole chain to the leaf observations,
so what an object depends on --- including the corporate actions behind
its frame --- is always legible. A derivation that \emph{fails} is not an
absence, and its treatment follows the split: a \textbf{projection}
failure --- a bootstrap with too few quotes --- is legible by
recomputation, never stored, and refreshes to success if a late arrival
fills the gap (MD-5); a \textbf{model-run} failure re-enters as a recorded
observation with its diagnostics, itself stale if its inputs later move; a
failed \emph{capture} (MD-2) is stored as an input-side arrival, not a
derivation. An absent object never passes for one never attempted. Record
enough --- observations, derived inputs, recipe, seed, lineage --- that
the object is a deterministic function of the record.
\trace{Reproducibility, Simulability, Correctness; C-2.2, C-3.2, C-14.9,
  C-14.11, C-14.15.}

\mlabel{MD-7}\textbf{MD-7. The prior is a term; the posterior is a
derived object.} A Bayesian calibration is not banned; it is placed. The
prior --- including any assumed dynamics for how a parameter moves
between observations --- is a declared, recorded term. The posterior is
a derived object of that prior-term, the recorded observations, and the
seed (MD-6): a projection where it computes over the record, a re-entered
observation where a filter evaluates a model. Either way it is \emph{not}
stored state --- no ``current posterior'' updates in place; today's
posterior is tomorrow's prior names a recomputation, not a mutable
variable. The framework mandates no dynamics --- a random walk, a
martingale, a mean reversion are recipe choices, never framework axioms
--- and whatever is declared is recorded, so the posterior is reproducible
on MD-6's terms.
\trace{Correctness, Reproducibility; C-2.8, C-4.11, C-14.15.}

\mlabel{MD-8}\textbf{MD-8. A broken state cannot hide.} A recurring
pathology of calibrated systems is the stored parameter that has drifted
from what the information implies: the system knows something it has not
yet written into its numbers. The Ledger defeats it by a mechanism that
differs by kind. A \textbf{projection} stores nothing --- recomputed from
the record on every read (C-4.11) --- so it cannot lag the record: the
drifted-parameter state is unrepresentable, and under a joint recipe an
on-demand read is the correct joint conditional expectation, so even two
correlated legs cannot hold a divergent pair. A \textbf{re-entered
observation} does store a number --- a fitted surface, or a joint
vol-and-correlation fit for a spread, sits on the record because the
ledger runs no model (C-14.9). It cannot drift in place: a new fit is a
new observation, forward, naming the inputs it consumed (MD-6), never an
edit of the old. And when an input it consumed is later corrected or
superseded --- a wrong price, or a corrected corporate action behind its
frame, both pinned lineage inputs (MD-10, MD-13) --- the ledger does not
silently recompute it; its lineage shows the input moved under it, and the
stale fit stands as a flagged open item for re-derivation, so the joint
spread surface cannot silently disagree with a corrected leg or a restated
split. That visibility is reached by default: a consumer reads a
re-entered observation through a projection --- ``the current fit for X'',
selecting the latest non-superseded observation and carrying its staleness
flag --- never the raw stored value. What the framework forbids everywhere is the
\emph{silent} divergence: for a projection it cannot arise; for a
re-entered observation it cannot hide.
\trace{Correctness; C-1.5, C-4.11, C-11.5, C-14.9.}

\mlabel{MD-9}\textbf{MD-9. Constraints live in the derived object;
no-arbitrage is declared, not assumed.} A recorded observation carries
no promise of mutual consistency --- a stale quote and a fresh one may
cross, and two sources may disagree; the record holds both. No-arbitrage,
smoothness, and monotonicity are constraints a derived object imposes on
its own output, not properties the record must have. A derived object
that cannot meet the constraints it declares records that it cannot ---
a visible diagnostic --- rather than serving an inadmissible object as
though it were sound. No-arbitrage is not negotiable for the object that
claims it, but it is enforced by detection, not admission: an
arbitrageable surface served as arbitrage-free is an economic error caught
after the fact by recomputation, not prevented at a door (C-13.2, C-13.3),
with claim and test both on the record. Which no-arbitrage conditions are
correct, and whether a fitted object hedges, are model choice, outside
this scope (C-Scope.11). No-arbitrage is the object's \emph{internal}
consistency; its \emph{external} validity --- whether it reprices what it
was drawn from --- is decided in price space (MD-15); both are recorded
diagnostics, neither prevented at a door.
\trace{Auditability, Correctness; C-13.2, C-13.3, C-Scope.11.}

\mlabel{MD-10}\textbf{MD-10. Corrections repair forward; the original is
never overwritten.} A market data fact found wrong is not edited. The
discovery is a new recorded observation that names the wrong one, and the
discrepancy stands as a visible open item. What happens downstream splits
by kind: every \textbf{projection} over the corrected fact recomputes
automatically on the next read; a \textbf{re-entered observation} that
consumed the wrong fact cannot re-run itself (C-14.9), so its recorded
lineage marks it stale and it stands as a flagged open item until a fresh
fit supersedes it forward (MD-8). A bulk retraction --- a vendor voiding
a whole day --- is many corrections at once; the explain may summarise
them as one named item rather than flood the book, so the signal is not
lost in the flags. Where money already moved on the strength of the wrong
fact, it moves back only as an authorised compensating transaction
(C-12.4). The observed value stays on the record exactly as observed.
Adjustments a later event forces on earlier data --- a split rescaling a
price, a dividend rescaling a cash-per-share figure --- are a \emph{change
of frame}, not a correction: computed at each read by the market data
operator (C-9.2, C-9.3), declared once per data-and-event pairing and
never improvised, obeying the algebra of MD-13. The original is preserved,
the adjusted value derived at read, and derived quantities are recomputed
from operator-adjusted inputs (C-9.3), never scalar-transported.
\trace{Auditability, Time travel, Correctness; C-9.3, C-12.4, C-12.6, C-14.9.}

\mlabel{MD-11}\textbf{MD-11. Simulated market data is real market data
under a different seed.} The state of the market data universe is the
initial state plus the fold of all recorded observations (C-2.8). A
stress scenario, a backtest, a Monte-Carlo path is not a separate system
but the same machine folding a generated stream of observations from a
recorded seed; production is simply the one path whose observations
happened to be real. Because the seed is the single non-record input and
it too is recorded, every simulated path replays exactly, on the same
terms as one over real observations --- but only because every model
output on the path is itself captured as a re-entered observation (MD-6);
a path that recorded only the seed could not replay.
\trace{Simulability, Reproducibility; C-2.8, C-10.1.}

\mlabel{MD-12}\textbf{MD-12. Derivation composes.} A derived object may
consume another derived object --- a volatility calibration reads a
discount curve, a spread model reads two surfaces. When it does, the
dependency --- which object, at which version and cut --- is part of the
recorded recipe (MD-6), so the composite is again a deterministic function
of the record and the MD-6 split composes: a chain of \textbf{projections}
reconstructs bit-for-bit from the record; a chain through a
\textbf{re-entered observation} enters that node as a \emph{leaf} --- read
back always, re-derived only with its retained model (MD-6). Because each
input is fixed at its own cut, never a mutable store, the ``mid-update
input'' cannot arise (MD-8): a curve feeding a calibration is consistent
with it by construction, not scheduling.
\trace{Reproducibility, Correctness; C-2.2, C-14.11.}

\mlabel{MD-13}\textbf{MD-13. A corporate action is a change of frame; the
operators that make it compose.} A corporate action is first-class --- a
recorded transaction like any other (C-9.1), never an adjustment applied
to market data outside the ledger. It re-coordinates every value recorded
before it (C-9.2), so a value is not a bare number but lives in a
\textbf{frame}: the coordinate system fixed by the corporate-action terms
in force \emph{and} the declared adjustment convention (which actions the
operator applies, price- or total-return, the rounding rule). The frame a
value arrives in is itself an \emph{asserted, recorded
fact} of the observation --- provenance registered per data kind (\S1),
never inferred from execution time, since a raw feed delivers unadjusted
but a vendor's back-adjusted series does not. The operator transports from
the \emph{declared} delivery frame, so a value already adjusted at source
is never adjusted again.

A corporate action carries its own three times. Its \textbf{ex-date} is
its as-of --- the moment the price re-coordinates, the frame boundary ---
distinct from when its terms became known. Before ex, or if the action is
cancelled, the operator is the identity on the price series, so a
cancelled-before-ex action never forces an un-adjust; an announcement, or
a revised dividend, is not a frame change but new information that refolds
projections (MD-5). And the operator is determined only once the terms are
\emph{resolved}, not merely announced: where the map turns on a later
fixing, proration, or holder election --- a scrip reference price, a
merger's blended consideration --- it exists from the \emph{resolution}
observation, and before that the frame is provisional and legible as such
(MD-6), never a silent wrong number.

The map is the Constitution's \textbf{market data operator} (C-9.2), a
projection computed at read, the original never overwritten (C-9.3,
MD-10). It acts on \emph{leaf} observations and is not in general a
proportional scalar --- a special cash dividend shifts a forward
additively, a listed option's OCC adjustment re-coordinates strike,
multiplier, and deliverable. Derived objects --- surfaces, curves,
correlations --- are never transported by an operator of their own;
\textbf{derived quantities are recomputed from operator-adjusted inputs}
(C-9.3, MD-6, MD-12). Operators compose at \emph{full precision}, the
minor-unit rounding (C-4.6) applied \emph{once} at read, so the composite
is grouping-independent and two parties reconstruct the identical value;
rounding between operators would not.

These operators form a plain algebra over frames. They \textbf{compose}:
two splits a week apart chain forward into one map, associative because it
is function composition. Over an action-free span the operator is the
\textbf{identity}, and some actions are the identity for a data kind --- a
proportional dividend figure stands across a split (C-9.2). An
\textbf{inverse} exists only where the transform preserves information;
under rounding most operators lose it and have none, and the framework
needs none --- the original is preserved (C-9.3), so any frame is reached
forward from the execution-time frame, never by inverting. The operator
changes a value's frame and leaves its as-of and as-at (MD-4) unchanged,
so it \emph{commutes} with as-of selection under a fixed as-at cut ---
adjust Tuesday into today's frame then pick Tuesday, or pick then adjust,
gives the same value. The cut must be fixed:
an action whose door time falls between two cuts is known only at the
later, so the composite operator differs between them. A corrected or late
corporate action is a \emph{superseding} event (MD-5, MD-10): every
projection over the old frame recomputes, and --- the action being a
pinned lineage input (MD-6) --- every re-entered observation whose lineage
reaches it is flagged stale, exactly as a corrected price does. Hence
\textbf{a quote is meaningful only given a time coordinate and a frame}:
any two parties who agree on (as-of, as-at, frame) reconstruct the
identical value.
\trace{Order, Time travel, Reproducibility, Correctness; C-9.1, C-9.2,
  C-9.3, C-2.3, C-4.6.}

\mlabel{MD-14}\textbf{MD-14. Every datum is dispute-ready.} Auditability
and reproducibility, pressed to their adversarial conclusion, give a
seventh governing principle: a dispute about a recorded value is settled
by \emph{replay}, not argument --- dispute-readiness is what those two
principles already are once an adversary insists, asking nothing new of
the record. To contest a mark is to contest a number, and the record
answers by exhibiting the whole chain that produced it --- the datum, its
provenance and attestation, the \textbf{cut} it was struck at, the
\textbf{frame} it was read in (MD-13), and the model it is bound to
(MD-15) --- and reproducing the valuation bit-for-bit (MD-6), the as-at
pinned to the mark's own cut, so the number that stood then is reproduced,
not a post-correction one. Its bound is exact: replay settles whether the
value is the faithful, deterministic consequence of the inputs the record
held at its stated (as-of, as-at, frame) --- the only question a system of
record can answer. It does \emph{not} settle a two-sided economic dispute:
a counterparty's mark, from its own surface, model, and frame, replays
bit-for-bit too, and which is right is model choice, out of scope (\S4).
But exhibiting the frame and the bound model \emph{localises} the
disagreement --- one side applied the split to the strike, the other to
the multiplier --- which is the value replay does deliver. Its reach is MD-6's:
a value reproduces from recorded prices unconditionally, a model's number
once the model is supplied (C-14.15), never wider.
\trace{Dispute-readiness (Auditability + Reproducibility); C-2.1, C-2.2,
  C-2.3, C-Scope.11.}

\mlabel{MD-15}\textbf{MD-15. A datum used in valuation is bound to a
model; validity is decided in price space.} A market data value on its
own is a number; it becomes a fact of the valuation universe only when
bound to the model through which it prices something. The binding is not a
claim that the model is true --- this manifesto asserts no model, and a
dynamics such as the martingale, or a latent ``true price'', is
admissible only as a declared, recorded term, never an axiom. The binding
is \textbf{recorded lineage}: which model, which version, stands beside
the datum (MD-6), exactly as an input cut does.

Why price space: no one argues a printed number is wrong in itself.
Disputes are about the valuation of \emph{units} (MD-14), and a unit's
value is defined only in price space, where truth is contractually fixed.
So a datum is validated where it is used: it is \textbf{valid on its
calibration set} when, through its bound model, it reprices the
instruments it was drawn from within a declared tolerance --- off-set use
is model extrapolation, whose validity is model choice, not this
round-trip. The round-trip is the acceptance test, and it shows
calibration and validation to be one act seen from opposite directions ---
calibration reads parameters from prices, validation reads prices back. Both are derived
objects over recorded inputs and the declared model term (MD-6), so
model-binding never undoes the split, and the repricing residual is
itself a re-entered observation --- a recorded diagnostic (MD-9), stale
if its inputs later move (MD-8), never a silent pass. Which model, and
which tolerance, are model choice, out of scope (C-Scope.11); that
validity is decided in price space is not.
\trace{Correctness; C-2.4, C-13.2, C-14.15, C-Scope.11.}

\section{3. The Life of an Observation}

One observation can walk the whole picture. Follow AAPL's official close
on a Tuesday, needed by an equity volatility surface.

\begin{itemize}
\tightlist
\item \textbf{Recorded once, in a frame.} The exchange publishes AAPL's
  Tuesday close of 150.20; the Transaction Executor admits a moveless
  observation-recording transaction. The value is on the record once
  (a redelivery is absorbed), stamped with execution time (Tue 16:00,
  as-of), monitor, and door time (Tue 18:30, as-at), in the pre-split
  frame it was quoted in (MD-1, MD-4, MD-13).
\item \textbf{A split is a change of frame.} The next week a two-for-one
  split is recorded as a first-class event. The market data operator
  carries 150.20 into the post-split frame as 75.10, computed at read; the
  original is never overwritten, and any later action's operator chains
  onto the split's, forward, in execution order (MD-13, MD-10).
\item \textbf{Corrected forward, two answers.} On Thursday the exchange
  restates the close to 150.50 --- a correction (a new observation naming
  the old, MD-10), distinct from the split. The two effects separate: the
  split alone reframes 150.20 to 75.10; the correction alone lifts 150.20
  to 150.50. So \emph{Tuesday as published Tuesday} (cut as-at Tue 18:30,
  pre-split frame) is 150.20, while \emph{Tuesday as valued now} (cut
  as-at today, post-split frame) is the corrected close reframed, 75.25 ---
  a quote needs both a time coordinate and a frame.
\item \textbf{Disputed.} A counterparty contests a collateral mark that
  used Tuesday's surface. It is not argued: replay exhibits 150.20, its
  provenance, cut, frame, and bound model, and reprices the mark
  bit-for-bit, the as-at pinned to the mark's cut (MD-14). Replay cannot
  decide whose surface is right, but it \emph{localises} the disagreement
  to a differing frame, input, or model.
\end{itemize}

\section{4. What This Manifesto Does Not Govern}

This manifesto governs how market data is recorded, ordered, corrected,
and projected. It does not govern \textbf{price formation and discovery}
(where a price comes from is not a question the record answers);
\textbf{which model, prior, or no-arbitrage condition is correct} (model
choice, whose test is hedging performance); \textbf{model governance,
validation, and retention} (a value reproduces from recorded prices
unconditionally, a model's number only if the model is kept); \textbf{the
terms of a corporate action} (the operators transport data through it,
MD-13; the terms arrive as recorded events, not decided here); and
\textbf{settlement} (the record projects committed activity toward it but
does not perform it). These lie outside the Ledger boundary and are
reconciled against it only at its perimeter (C-Scope.11).

\section{5. Subordination}

This manifesto is subordinate to the Constitution and rated against it,
never the other way round. Its vocabulary is the Constitution's, one name
per component, no synonym for any fixed term. Where a principle here would
conflict with the Constitution, the Constitution prevails and the
conflict is parked with the exact amendment text and the exact clause it
replaces --- never resolved by quiet narrowing. A genuine departure amends
this document explicitly first; an article identifier is never reused or
renumbered.

\section{Conclusion}

Reconciliation of market data has the same root cause as every other: the
same fact stored more than once, and the Ledger removes it. An observation
is recorded once; every curve, surface, and calibration is a deterministic
function of what the record holds --- read back always, re-derived only
where a model runs; every value carries its times and its frame, is bound
to a model, and is settled by replay; corrections repair forward,
originals survive. Operate it any other way and you reintroduce the
divergence the architecture exists to remove; this way, there is nothing
left to reconcile.

\noindent\textbf{Amendment record --- Version 1.1} embeds three mandates
in the certified 1.0: corporate-action frames and the operator algebra
(MD-13); dispute-readiness, a derived seventh principle (MD-14); and model
binding with price-space validation (MD-15). MD-4, MD-6, MD-9, MD-10
extended; MD-1--MD-12 keep their numbers.

\end{document}
